% Options for packages loaded elsewhere
\PassOptionsToPackage{unicode}{hyperref}
\PassOptionsToPackage{hyphens}{url}
%
\documentclass[
]{article}
\usepackage{amsmath,amssymb}
\usepackage{iftex}
\ifPDFTeX
  \usepackage[T1]{fontenc}
  \usepackage[utf8]{inputenc}
  \usepackage{textcomp} % provide euro and other symbols
\else % if luatex or xetex
  \usepackage{unicode-math} % this also loads fontspec
  \defaultfontfeatures{Scale=MatchLowercase}
  \defaultfontfeatures[\rmfamily]{Ligatures=TeX,Scale=1}
\fi
\usepackage{lmodern}
\ifPDFTeX\else
  % xetex/luatex font selection
\fi
% Use upquote if available, for straight quotes in verbatim environments
\IfFileExists{upquote.sty}{\usepackage{upquote}}{}
\IfFileExists{microtype.sty}{% use microtype if available
  \usepackage[]{microtype}
  \UseMicrotypeSet[protrusion]{basicmath} % disable protrusion for tt fonts
}{}
\makeatletter
\@ifundefined{KOMAClassName}{% if non-KOMA class
  \IfFileExists{parskip.sty}{%
    \usepackage{parskip}
  }{% else
    \setlength{\parindent}{0pt}
    \setlength{\parskip}{6pt plus 2pt minus 1pt}}
}{% if KOMA class
  \KOMAoptions{parskip=half}}
\makeatother
\usepackage{xcolor}
\usepackage[margin=1in]{geometry}
\usepackage{color}
\usepackage{fancyvrb}
\newcommand{\VerbBar}{|}
\newcommand{\VERB}{\Verb[commandchars=\\\{\}]}
\DefineVerbatimEnvironment{Highlighting}{Verbatim}{commandchars=\\\{\}}
% Add ',fontsize=\small' for more characters per line
\usepackage{framed}
\definecolor{shadecolor}{RGB}{248,248,248}
\newenvironment{Shaded}{\begin{snugshade}}{\end{snugshade}}
\newcommand{\AlertTok}[1]{\textcolor[rgb]{0.94,0.16,0.16}{#1}}
\newcommand{\AnnotationTok}[1]{\textcolor[rgb]{0.56,0.35,0.01}{\textbf{\textit{#1}}}}
\newcommand{\AttributeTok}[1]{\textcolor[rgb]{0.13,0.29,0.53}{#1}}
\newcommand{\BaseNTok}[1]{\textcolor[rgb]{0.00,0.00,0.81}{#1}}
\newcommand{\BuiltInTok}[1]{#1}
\newcommand{\CharTok}[1]{\textcolor[rgb]{0.31,0.60,0.02}{#1}}
\newcommand{\CommentTok}[1]{\textcolor[rgb]{0.56,0.35,0.01}{\textit{#1}}}
\newcommand{\CommentVarTok}[1]{\textcolor[rgb]{0.56,0.35,0.01}{\textbf{\textit{#1}}}}
\newcommand{\ConstantTok}[1]{\textcolor[rgb]{0.56,0.35,0.01}{#1}}
\newcommand{\ControlFlowTok}[1]{\textcolor[rgb]{0.13,0.29,0.53}{\textbf{#1}}}
\newcommand{\DataTypeTok}[1]{\textcolor[rgb]{0.13,0.29,0.53}{#1}}
\newcommand{\DecValTok}[1]{\textcolor[rgb]{0.00,0.00,0.81}{#1}}
\newcommand{\DocumentationTok}[1]{\textcolor[rgb]{0.56,0.35,0.01}{\textbf{\textit{#1}}}}
\newcommand{\ErrorTok}[1]{\textcolor[rgb]{0.64,0.00,0.00}{\textbf{#1}}}
\newcommand{\ExtensionTok}[1]{#1}
\newcommand{\FloatTok}[1]{\textcolor[rgb]{0.00,0.00,0.81}{#1}}
\newcommand{\FunctionTok}[1]{\textcolor[rgb]{0.13,0.29,0.53}{\textbf{#1}}}
\newcommand{\ImportTok}[1]{#1}
\newcommand{\InformationTok}[1]{\textcolor[rgb]{0.56,0.35,0.01}{\textbf{\textit{#1}}}}
\newcommand{\KeywordTok}[1]{\textcolor[rgb]{0.13,0.29,0.53}{\textbf{#1}}}
\newcommand{\NormalTok}[1]{#1}
\newcommand{\OperatorTok}[1]{\textcolor[rgb]{0.81,0.36,0.00}{\textbf{#1}}}
\newcommand{\OtherTok}[1]{\textcolor[rgb]{0.56,0.35,0.01}{#1}}
\newcommand{\PreprocessorTok}[1]{\textcolor[rgb]{0.56,0.35,0.01}{\textit{#1}}}
\newcommand{\RegionMarkerTok}[1]{#1}
\newcommand{\SpecialCharTok}[1]{\textcolor[rgb]{0.81,0.36,0.00}{\textbf{#1}}}
\newcommand{\SpecialStringTok}[1]{\textcolor[rgb]{0.31,0.60,0.02}{#1}}
\newcommand{\StringTok}[1]{\textcolor[rgb]{0.31,0.60,0.02}{#1}}
\newcommand{\VariableTok}[1]{\textcolor[rgb]{0.00,0.00,0.00}{#1}}
\newcommand{\VerbatimStringTok}[1]{\textcolor[rgb]{0.31,0.60,0.02}{#1}}
\newcommand{\WarningTok}[1]{\textcolor[rgb]{0.56,0.35,0.01}{\textbf{\textit{#1}}}}
\usepackage{graphicx}
\makeatletter
\def\maxwidth{\ifdim\Gin@nat@width>\linewidth\linewidth\else\Gin@nat@width\fi}
\def\maxheight{\ifdim\Gin@nat@height>\textheight\textheight\else\Gin@nat@height\fi}
\makeatother
% Scale images if necessary, so that they will not overflow the page
% margins by default, and it is still possible to overwrite the defaults
% using explicit options in \includegraphics[width, height, ...]{}
\setkeys{Gin}{width=\maxwidth,height=\maxheight,keepaspectratio}
% Set default figure placement to htbp
\makeatletter
\def\fps@figure{htbp}
\makeatother
\setlength{\emergencystretch}{3em} % prevent overfull lines
\providecommand{\tightlist}{%
  \setlength{\itemsep}{0pt}\setlength{\parskip}{0pt}}
\setcounter{secnumdepth}{5}
\usepackage{booktabs}
\usepackage{caption}
\usepackage{longtable}
\usepackage{colortbl}
\usepackage{array}
\usepackage{anyfontsize}
\usepackage{multirow}
\ifLuaTeX
  \usepackage{selnolig}  % disable illegal ligatures
\fi
\usepackage{bookmark}
\IfFileExists{xurl.sty}{\usepackage{xurl}}{} % add URL line breaks if available
\urlstyle{same}
\hypersetup{
  pdftitle={Analyse des données ménages - Sect1 Harvest W4},
  pdfauthor={DIALLO Mamadou Moussa - DIOP Papa Boubacar},
  hidelinks,
  pdfcreator={LaTeX via pandoc}}

\title{Analyse des données ménages - Sect1 Harvest W4}
\author{DIALLO Mamadou Moussa - DIOP Papa Boubacar}
\date{2026-03-17}

\begin{document}
\maketitle

{
\setcounter{tocdepth}{2}
\tableofcontents
}
L'objectif de ce tp est de décrire la composition et les
caractéristiques démographiques des ménages enquêtés dans les quatres
vagues du GHS Panel à partir de la section sect1\_harvest.

\begin{Shaded}
\begin{Highlighting}[]
\CommentTok{\# Importation des packages}
\FunctionTok{library}\NormalTok{(haven)}
\FunctionTok{library}\NormalTok{(dplyr)}
\FunctionTok{library}\NormalTok{(ggplot2)}
\FunctionTok{library}\NormalTok{(apyramid)}
\FunctionTok{library}\NormalTok{(naniar)}
\FunctionTok{library}\NormalTok{(gtsummary)}
\FunctionTok{library}\NormalTok{(rstatix)}
\FunctionTok{library}\NormalTok{(PropCIs)}
\FunctionTok{library}\NormalTok{(patchwork)}
\FunctionTok{library}\NormalTok{(tidyr)}
\FunctionTok{library}\NormalTok{(purrr)}
\end{Highlighting}
\end{Shaded}

\begin{Shaded}
\begin{Highlighting}[]
\CommentTok{\# Chargement des données}
\NormalTok{data }\OtherTok{\textless{}{-}} \FunctionTok{read\_stata}\NormalTok{(}\StringTok{"../data/raw/sect1\_harvestw4.dta"}\NormalTok{)}
\CommentTok{\# data\_2 \textless{}{-} read\_stata("../data/raw/secta\_harvestw4.dta")}
\end{Highlighting}
\end{Shaded}

\section{Exploration initiale des
données}\label{exploration-initiale-des-donnuxe9es}

\subsection{Structure des données}\label{structure-des-donnuxe9es}

\subsection{Vérification des
doublons}\label{vuxe9rification-des-doublons}

\begin{Shaded}
\begin{Highlighting}[]
\NormalTok{data }\SpecialCharTok{\%\textgreater{}\%} \FunctionTok{count}\NormalTok{(hhid, indiv) }\SpecialCharTok{\%\textgreater{}\%} \FunctionTok{filter}\NormalTok{(n }\SpecialCharTok{\textgreater{}} \DecValTok{1}\NormalTok{) }\SpecialCharTok{\%\textgreater{}\%} \FunctionTok{count}\NormalTok{()}
\end{Highlighting}
\end{Shaded}

\begin{verbatim}
## # A tibble: 1 x 1
##       n
##   <int>
## 1     0
\end{verbatim}

Il n'y a donc pas de doublons d'individu dans la base sect1\_harvest.

\subsection{Valeurs manquantes}\label{valeurs-manquantes}

\begin{Shaded}
\begin{Highlighting}[]
\FunctionTok{vis\_miss}\NormalTok{(data, }\AttributeTok{warn\_large\_data =}\NormalTok{ F)}
\end{Highlighting}
\end{Shaded}

\includegraphics{rapport_final-Rmd_files/figure-latex/missing-values-1.pdf}

Le graphique ci-dessus représente le schéma des valeurs manquantes nous
montre qu'il y a 59,1\% de valeurs manquantes. Cependant, cette valeur
est à prendre avec des pincettes car il y a aussi bien de vraies valeurs
manquantes que de fausses valeurs manquantes. Ces fausses valeurs
manquantes correspondent à des sauts dans le questionnaire.

\section{Analyse univariée de l'âge des
membres}\label{analyse-univariuxe9e-de-luxe2ge-des-membres}

\subsection{Visualisations}\label{visualisations}

\begin{Shaded}
\begin{Highlighting}[]
\CommentTok{\# Histogramme avec des classes de 5 ans}
\FunctionTok{ggplot}\NormalTok{(data, }\FunctionTok{aes}\NormalTok{(}\AttributeTok{x =}\NormalTok{ s1q4)) }\SpecialCharTok{+}
  \FunctionTok{geom\_histogram}\NormalTok{(}\AttributeTok{binwidth =} \DecValTok{5}\NormalTok{, }\AttributeTok{fill =} \StringTok{"darkgreen"}\NormalTok{, }\AttributeTok{color =} \StringTok{"black"}\NormalTok{) }\SpecialCharTok{+}
  \FunctionTok{labs}\NormalTok{(}
    \AttributeTok{title =} \StringTok{"Distribution de l\textquotesingle{}âge des membres"}\NormalTok{,}
    \AttributeTok{x =} \StringTok{"Âge (années)"}\NormalTok{,}
    \AttributeTok{y =} \StringTok{"Effectif"}
\NormalTok{  ) }\SpecialCharTok{+}
  \FunctionTok{theme\_minimal}\NormalTok{()}
\end{Highlighting}
\end{Shaded}

\includegraphics{rapport_final-Rmd_files/figure-latex/age-histogramme-1.pdf}
L'histogramme ci-dessus montre la répartition des âges des membres des
ménages par tranche de 5 ans. On peut observer que certaines tranches
sont plus représentées, ce qui permet de voir si la population est
plutôt jeune ou âgée. On constate que la distribution est étalée à
droite, ce qui montre que la population est relativement jeune.

\begin{Shaded}
\begin{Highlighting}[]
\CommentTok{\# Boîte à moustaches}
\FunctionTok{ggplot}\NormalTok{(data, }\FunctionTok{aes}\NormalTok{(}\AttributeTok{x =}\NormalTok{ s1q4)) }\SpecialCharTok{+}
  \FunctionTok{geom\_boxplot}\NormalTok{(}\AttributeTok{fill =} \StringTok{"lightblue"}\NormalTok{) }\SpecialCharTok{+}
  \FunctionTok{labs}\NormalTok{(}
    \AttributeTok{title =} \StringTok{"Distribution de l\textquotesingle{}âge"}\NormalTok{,}
    \AttributeTok{x =} \StringTok{"Âge (années)"}
\NormalTok{  ) }\SpecialCharTok{+}
  \FunctionTok{theme\_minimal}\NormalTok{()}
\end{Highlighting}
\end{Shaded}

\includegraphics{rapport_final-Rmd_files/figure-latex/age-boxplot-1.pdf}

L'analyse de la boîte à moustaches renforce le constat de la jeunesse.
En effet, 75\% des individus ont un âge compris entre 8 et 36ans.
Cependant, nous constatons aussi la présence de valeurs extrêmes,
c'est-à-dire les valeurs qui dépassent la moustache droite.

\subsection{Statistiques descriptives}\label{statistiques-descriptives}

\begin{Shaded}
\begin{Highlighting}[]
\NormalTok{stat\_desc\_age }\OtherTok{\textless{}{-}}\NormalTok{ data }\SpecialCharTok{\%\textgreater{}\%} 
  \FunctionTok{select}\NormalTok{(s1q4) }\SpecialCharTok{\%\textgreater{}\%}
  \FunctionTok{tbl\_summary}\NormalTok{(}
    \AttributeTok{type =}\NormalTok{ s1q4 }\SpecialCharTok{\textasciitilde{}} \StringTok{"continuous2"}\NormalTok{,}
    \AttributeTok{statistic =} \FunctionTok{all\_continuous}\NormalTok{() }\SpecialCharTok{\textasciitilde{}} \FunctionTok{c}\NormalTok{(}\StringTok{"\{mean\}"}\NormalTok{, }\StringTok{"\{median\}"}\NormalTok{, }\StringTok{"\{p25\}"}\NormalTok{, }\StringTok{"\{p75\}"}\NormalTok{),}
    \AttributeTok{label =}\NormalTok{ s1q4 }\SpecialCharTok{\textasciitilde{}} \StringTok{"Âge"}
\NormalTok{  )}
\NormalTok{stat\_desc\_age}
\end{Highlighting}
\end{Shaded}

\begin{table}[t]
\fontsize{12.0pt}{14.0pt}\selectfont
\begin{tabular*}{\linewidth}{@{\extracolsep{\fill}}lc}
\toprule
\textbf{Characteristic} & \textbf{N = 30,337} \\ 
\midrule\addlinespace[2.5pt]
Âge &  \\ 
    Mean & 24 \\ 
    Median & 18 \\ 
    Q1 & 8 \\ 
    Q3 & 36 \\ 
    Unknown & 3,780 \\ 
\bottomrule
\end{tabular*}
\end{table}

L'âge moyen des individus de l'échantillon est de 24 ans. Cette
statistique étant sensible aux valeurs extrêmes, l'âge médian constitue
une alternative plus robuste à ces valeurs extrêmes. L'âge médian dans
l'échantillon est de 18 ans. Et seulement 25\% ont un âge supérieur à
36ans. Ainsi, ces statistiques montrent que la population étudiée est
jeune.

\subsection{Test de normalité}\label{test-de-normalituxe9}

\begin{Shaded}
\begin{Highlighting}[]
\CommentTok{\# Test de Shapiro{-}Wilk (taille de l\textquotesingle{}échantillon est supérieur à 5)}
\CommentTok{\# shapiro.test(data$s1q4)}
\FunctionTok{library}\NormalTok{(tseries)}
\FunctionTok{jarque.bera.test}\NormalTok{(}\FunctionTok{na.omit}\NormalTok{(data}\SpecialCharTok{$}\NormalTok{s1q4))}
\end{Highlighting}
\end{Shaded}

\begin{verbatim}
## 
##  Jarque Bera Test
## 
## data:  na.omit(data$s1q4)
## X-squared = 4503.7, df = 2, p-value < 2.2e-16
\end{verbatim}

\begin{Shaded}
\begin{Highlighting}[]
\FunctionTok{ggplot}\NormalTok{(}\AttributeTok{data =} \FunctionTok{data.frame}\NormalTok{(}\AttributeTok{s1q4\_clean =} \FunctionTok{na.omit}\NormalTok{(data}\SpecialCharTok{$}\NormalTok{s1q4)), }\FunctionTok{aes}\NormalTok{(}\AttributeTok{sample =} \FunctionTok{na.omit}\NormalTok{(data}\SpecialCharTok{$}\NormalTok{s1q4))) }\SpecialCharTok{+}
  \FunctionTok{stat\_qq}\NormalTok{() }\SpecialCharTok{+}
  \FunctionTok{stat\_qq\_line}\NormalTok{(}\AttributeTok{col =} \StringTok{"red"}\NormalTok{)}
\end{Highlighting}
\end{Shaded}

\includegraphics{rapport_final-Rmd_files/figure-latex/age-normalite-1.pdf}
L'analyse visuelle de l'histogramme des âges montre que l'âge ne semble
pas normalement distribué. Pour confirmer ce constat, on peut utiliser
un test statistique de normalité. Cependant, la taille de l'échantillon
étant supérieure à 5000, le test de shpiro-wilk ne peut s'effectuer.
Néanmoins, une p-value très faible pour le test de normalité de Jarque
Bera nous permet de conclure que la variable âge n'est pas normalement
distribuée. De plus, le graphique quantile-quantile permet de vérifier
que la normalité n'est pas respectée.

\section{Pyramide des âges}\label{pyramide-des-uxe2ges}

\begin{Shaded}
\begin{Highlighting}[]
\NormalTok{data\_pyramid }\OtherTok{\textless{}{-}}\NormalTok{ data }\SpecialCharTok{\%\textgreater{}\%} 
  \FunctionTok{select}\NormalTok{(}\AttributeTok{age =}\NormalTok{ s1q4, }\AttributeTok{sexe =}\NormalTok{ s1q2) }\SpecialCharTok{\%\textgreater{}\%} 
  \FunctionTok{filter}\NormalTok{(}\SpecialCharTok{!}\FunctionTok{is.na}\NormalTok{(age), }\SpecialCharTok{!}\FunctionTok{is.na}\NormalTok{(sexe)) }\SpecialCharTok{\%\textgreater{}\%} 
  \FunctionTok{mutate}\NormalTok{(}\AttributeTok{age\_group =} \FunctionTok{cut}\NormalTok{(age, }\AttributeTok{breaks =} \FunctionTok{seq}\NormalTok{(}\DecValTok{0}\NormalTok{, }\DecValTok{100}\NormalTok{, }\AttributeTok{by =} \DecValTok{5}\NormalTok{)))}

\FunctionTok{age\_pyramid}\NormalTok{(data\_pyramid, }\AttributeTok{age\_group =}\NormalTok{ age\_group, }\AttributeTok{split\_by =}\NormalTok{ sexe) }\SpecialCharTok{+}
  \FunctionTok{labs}\NormalTok{(}\AttributeTok{title =} \StringTok{"Pyramide des âges"}\NormalTok{) }\SpecialCharTok{+}
  \FunctionTok{theme\_minimal}\NormalTok{()}
\end{Highlighting}
\end{Shaded}

\includegraphics{rapport_final-Rmd_files/figure-latex/pyramide-1.pdf}

La pyramide des âges montre la répartition de la population par âge et
par sexe. On constate que le pyramide a une base élargie et un sommet
pointu et les effectifs diminuent avec l'âge. Ce qui confirme la
jeunesse de la population. Cependant, on constate que l'effectif des
individus de 5 à 9 ans est supérieur à celui des enfants de 0 à 4 ans.

\section{Fréquence des liens de
parenté}\label{fruxe9quence-des-liens-de-parentuxe9}

\subsection{Diagramme en barres
simple}\label{diagramme-en-barres-simple}

\begin{Shaded}
\begin{Highlighting}[]
\FunctionTok{library}\NormalTok{(forcats)}
\NormalTok{data }\SpecialCharTok{\%\textgreater{}\%} 
  \FunctionTok{count}\NormalTok{(s1q3) }\SpecialCharTok{\%\textgreater{}\%}
  \FunctionTok{arrange}\NormalTok{(}\FunctionTok{desc}\NormalTok{(n)) }\SpecialCharTok{\%\textgreater{}\%} 
  \FunctionTok{ggplot}\NormalTok{(}\FunctionTok{aes}\NormalTok{(}\AttributeTok{x =}\NormalTok{ n, }\AttributeTok{y =} \FunctionTok{fct\_reorder}\NormalTok{(}\FunctionTok{as\_factor}\NormalTok{(s1q3), n))) }\SpecialCharTok{+}  \CommentTok{\# conversion temporaire}
  \FunctionTok{geom\_col}\NormalTok{(}\AttributeTok{fill =} \StringTok{"darkblue"}\NormalTok{) }\SpecialCharTok{+}
  \FunctionTok{labs}\NormalTok{(}
    \AttributeTok{title =} \StringTok{"Fréquence des liens de parenté"}\NormalTok{,}
    \AttributeTok{x =} \StringTok{"Fréquence"}\NormalTok{,}
    \AttributeTok{y =} \StringTok{""}
\NormalTok{  ) }\SpecialCharTok{+}
  \FunctionTok{theme\_minimal}\NormalTok{()}
\end{Highlighting}
\end{Shaded}

\includegraphics{rapport_final-Rmd_files/figure-latex/parente-barres-1.pdf}

Ce graphique ci-dessus nous montre les effectifs par lien de parenté
avec le chef de ménage des différents individus de l'échantillon. On
constate que les individus qui sont enfants uniques du chef de ménage
représente la part la plus importante. A l'opposé, les individus qui
sont de la belle famille comptent parmi les effectifs les plus faibles.

\subsection{Proportions avec intervalles de
confiance}\label{proportions-avec-intervalles-de-confiance}

\begin{Shaded}
\begin{Highlighting}[]
\NormalTok{total }\OtherTok{\textless{}{-}} \FunctionTok{nrow}\NormalTok{(data)}

\NormalTok{resultats\_parente }\OtherTok{\textless{}{-}}\NormalTok{ data }\SpecialCharTok{\%\textgreater{}\%}
  \FunctionTok{filter}\NormalTok{(}\SpecialCharTok{!}\FunctionTok{is.na}\NormalTok{(s1q3)) }\SpecialCharTok{\%\textgreater{}\%}
  \FunctionTok{count}\NormalTok{(s1q3, }\AttributeTok{name =} \StringTok{"effectif"}\NormalTok{) }\SpecialCharTok{\%\textgreater{}\%}
  \FunctionTok{mutate}\NormalTok{(}
    \AttributeTok{proportion =}\NormalTok{ effectif }\SpecialCharTok{/}\NormalTok{ total,}
    \AttributeTok{proportion\_pct =}\NormalTok{ scales}\SpecialCharTok{::}\FunctionTok{percent}\NormalTok{(proportion, }\AttributeTok{accuracy =} \FloatTok{0.1}\NormalTok{),}
    \AttributeTok{IC\_lower =} \FunctionTok{sapply}\NormalTok{(effectif, }\ControlFlowTok{function}\NormalTok{(x) }\FunctionTok{scoreci}\NormalTok{(x, total, }\AttributeTok{conf.level =} \FloatTok{0.95}\NormalTok{)}\SpecialCharTok{$}\NormalTok{conf.int[}\DecValTok{1}\NormalTok{]),}
    \AttributeTok{IC\_upper =} \FunctionTok{sapply}\NormalTok{(effectif, }\ControlFlowTok{function}\NormalTok{(x) }\FunctionTok{scoreci}\NormalTok{(x, total, }\AttributeTok{conf.level =} \FloatTok{0.95}\NormalTok{)}\SpecialCharTok{$}\NormalTok{conf.int[}\DecValTok{2}\NormalTok{]),}
    \AttributeTok{IC\_text =} \FunctionTok{paste0}\NormalTok{(}\StringTok{"["}\NormalTok{, scales}\SpecialCharTok{::}\FunctionTok{percent}\NormalTok{(IC\_lower, }\AttributeTok{accuracy =} \FloatTok{0.1}\NormalTok{), }\StringTok{" ; "}\NormalTok{, }
\NormalTok{                     scales}\SpecialCharTok{::}\FunctionTok{percent}\NormalTok{(IC\_upper, }\AttributeTok{accuracy =} \FloatTok{0.1}\NormalTok{), }\StringTok{"]"}\NormalTok{)}
\NormalTok{  )}

\CommentTok{\# Affichage du tableau}
\NormalTok{resultats\_parente }\SpecialCharTok{\%\textgreater{}\%}
  \FunctionTok{select}\NormalTok{(s1q3, effectif, proportion\_pct, IC\_text) }\SpecialCharTok{\%\textgreater{}\%}
  \FunctionTok{rename}\NormalTok{(}
    \StringTok{\textasciigrave{}}\AttributeTok{Lien de parenté}\StringTok{\textasciigrave{}} \OtherTok{=}\NormalTok{ (s1q3),}
    \AttributeTok{Effectif =}\NormalTok{ effectif,}
    \AttributeTok{Proportion =}\NormalTok{ proportion\_pct,}
    \StringTok{\textasciigrave{}}\AttributeTok{IC 95\%}\StringTok{\textasciigrave{}} \OtherTok{=}\NormalTok{ IC\_text}
\NormalTok{  )}
\end{Highlighting}
\end{Shaded}

\begin{verbatim}
## # A tibble: 14 x 4
##    `Lien de parenté`                     Effectif Proportion `IC 95%`       
##    <dbl+lbl>                                <int> <chr>      <chr>          
##  1  1 [1. HEAD]                              4980 16.4%      [16.0% ; 16.8%]
##  2  2 [2. SPOUSE]                            4273 14.1%      [13.7% ; 14.5%]
##  3  3 [3. OWN CHILD]                        14636 48.2%      [47.7% ; 48.8%]
##  4  4 [4. STEP CHILD]                         192 0.6%       [0.6% ; 0.7%]  
##  5  5 [5. ADOPTED CHILD]                       88 0.3%       [0.2% ; 0.4%]  
##  6  6 [6. GRANDCHILD]                        1045 3.4%       [3.2% ; 3.7%]  
##  7  7 [7. BROTHER/SISTER]                     487 1.6%       [1.5% ; 1.8%]  
##  8  8 [8. NIECE/NEPHEW]                       229 0.8%       [0.7% ; 0.9%]  
##  9  9 [9. BROTHER/SISTER-IN-LAW]              137 0.5%       [0.4% ; 0.5%]  
## 10 10 [10. PARENT]                            291 1.0%       [0.9% ; 1.1%]  
## 11 11 [11. PARENT-IN-LAW]                      20 0.1%       [0.0% ; 0.1%]  
## 12 12 [12. DOMESTIC HELP]                      39 0.1%       [0.1% ; 0.2%]  
## 13 14 [14. OTHER RELATION (SPECIFY)]          123 0.4%       [0.3% ; 0.5%]  
## 14 15 [15. OTHER NON-RELATION (SPECIFY)]       17 0.1%       [0.0% ; 0.1%]
\end{verbatim}

\begin{Shaded}
\begin{Highlighting}[]
\CommentTok{\# Visualisation avec IC}
\NormalTok{resultats\_parente }\SpecialCharTok{\%\textgreater{}\%}
  \FunctionTok{ggplot}\NormalTok{(}\FunctionTok{aes}\NormalTok{(}\AttributeTok{x =}\NormalTok{ proportion, }\AttributeTok{y =} \FunctionTok{fct\_reorder}\NormalTok{(}\FunctionTok{as\_factor}\NormalTok{(s1q3), proportion))) }\SpecialCharTok{+}
  \FunctionTok{geom\_col}\NormalTok{(}\AttributeTok{fill =} \StringTok{"darkblue"}\NormalTok{) }\SpecialCharTok{+}
  \FunctionTok{geom\_errorbar}\NormalTok{(}\FunctionTok{aes}\NormalTok{(}\AttributeTok{xmin =}\NormalTok{ IC\_lower, }\AttributeTok{xmax =}\NormalTok{ IC\_upper), }\AttributeTok{width =} \FloatTok{0.2}\NormalTok{) }\SpecialCharTok{+}
  \FunctionTok{labs}\NormalTok{(}
    \AttributeTok{title =} \StringTok{"Proportions des liens de parenté avec IC à 95\%"}\NormalTok{,}
    \AttributeTok{x =} \StringTok{"Proportion"}\NormalTok{,}
    \AttributeTok{y =} \StringTok{""}
\NormalTok{  ) }\SpecialCharTok{+}
  \FunctionTok{scale\_x\_continuous}\NormalTok{(}\AttributeTok{labels =}\NormalTok{ scales}\SpecialCharTok{::}\NormalTok{percent) }\SpecialCharTok{+}
  \FunctionTok{theme\_minimal}\NormalTok{()}
\end{Highlighting}
\end{Shaded}

\includegraphics{rapport_final-Rmd_files/figure-latex/parente-proportions-ic-1.pdf}

Ce graphique représente les effectifs classés par lien de parenté avec
le chef de ménage avec un intervalle de confiance à 95\%.

\section{Comparaison de la taille des ménages entre zones rurales et
urbaines}\label{comparaison-de-la-taille-des-muxe9nages-entre-zones-rurales-et-urbaines}

\subsection{Préparation des
données}\label{pruxe9paration-des-donnuxe9es}

\begin{Shaded}
\begin{Highlighting}[]
\CommentTok{\# Calcul de la taille des ménages}
\NormalTok{data\_group }\OtherTok{\textless{}{-}}\NormalTok{ data }\SpecialCharTok{\%\textgreater{}\%}  
  \FunctionTok{select}\NormalTok{(hhid, sector) }\SpecialCharTok{\%\textgreater{}\%} 
  \FunctionTok{group\_by}\NormalTok{(hhid) }\SpecialCharTok{\%\textgreater{}\%} 
  \FunctionTok{mutate}\NormalTok{(}\AttributeTok{taille\_men =} \FunctionTok{n}\NormalTok{()) }\SpecialCharTok{\%\textgreater{}\%} \FunctionTok{ungroup}\NormalTok{()}
\end{Highlighting}
\end{Shaded}

\subsection{Visualisation}\label{visualisation}

\begin{Shaded}
\begin{Highlighting}[]
\FunctionTok{ggplot}\NormalTok{(data\_group, }\FunctionTok{aes}\NormalTok{(}\AttributeTok{x =} \FunctionTok{as\_factor}\NormalTok{(sector), }\AttributeTok{y =}\NormalTok{ taille\_men)) }\SpecialCharTok{+}
  \FunctionTok{geom\_boxplot}\NormalTok{(}\AttributeTok{fill =} \StringTok{"lightblue"}\NormalTok{) }\SpecialCharTok{+}
  \FunctionTok{labs}\NormalTok{(}
    \AttributeTok{title =} \StringTok{"Distribution de la taille des ménages par zone"}\NormalTok{,}
    \AttributeTok{x =} \StringTok{"Zone"}\NormalTok{,}
    \AttributeTok{y =} \StringTok{"Taille du ménage"}
\NormalTok{  ) }\SpecialCharTok{+}
  \FunctionTok{theme\_minimal}\NormalTok{()}
\end{Highlighting}
\end{Shaded}

\includegraphics{rapport_final-Rmd_files/figure-latex/boxplot-taille-zone-1.pdf}

Le graphique ci-dessus représente le boxplot de la taille du ménage
selon le milieu de résidence. On constate que la taille médiane du
ménage est inférieure en milieu urbain qu'en milieu rural. Cependant,
pour vérifier si cette différence est significative, on peut faire
recours à un test de moyenne entre les deux groupes. \#\# Test
statistique

\begin{Shaded}
\begin{Highlighting}[]
\FunctionTok{wilcox.test}\NormalTok{(taille\_men }\SpecialCharTok{\textasciitilde{}} \FunctionTok{as\_factor}\NormalTok{(sector), }\AttributeTok{data =}\NormalTok{ data\_group)}
\end{Highlighting}
\end{Shaded}

\begin{verbatim}
## 
##  Wilcoxon rank sum test with continuity correction
## 
## data:  taille_men by as_factor(sector)
## W = 76136889, p-value < 2.2e-16
## alternative hypothesis: true location shift is not equal to 0
\end{verbatim}

Le test de significativité de la différence de moyenne entre les deux
groupes nous donne une p-value très inférieure à 0,05. Ce qui nous
permet de conclure que la différence de moyenne entre les deux groupes
est significative.

\section{Tableau récapitulatif stratifié par
zone}\label{tableau-ruxe9capitulatif-stratifiuxe9-par-zone}

\subsection{Préparation des données
complètes}\label{pruxe9paration-des-donnuxe9es-compluxe8tes}

\begin{Shaded}
\begin{Highlighting}[]
\NormalTok{data }\OtherTok{\textless{}{-}}\NormalTok{ data }\SpecialCharTok{\%\textgreater{}\%} 
  \FunctionTok{group\_by}\NormalTok{(hhid) }\SpecialCharTok{\%\textgreater{}\%} 
  \FunctionTok{mutate}\NormalTok{(}\AttributeTok{taille\_men =} \FunctionTok{n}\NormalTok{()) }\SpecialCharTok{\%\textgreater{}\%} 
  \FunctionTok{ungroup}\NormalTok{()}

\NormalTok{data\_tab }\OtherTok{\textless{}{-}}\NormalTok{ data }\SpecialCharTok{\%\textgreater{}\%}
  \FunctionTok{select}\NormalTok{(hhid, s1q4, s1q2, taille\_men, sector) }\SpecialCharTok{\%\textgreater{}\%}
  \FunctionTok{mutate}\NormalTok{(}
    \AttributeTok{sector =} \FunctionTok{as\_factor}\NormalTok{(sector),}
    \AttributeTok{s1q2 =} \FunctionTok{as\_factor}\NormalTok{(s1q2)}
\NormalTok{  )}
\end{Highlighting}
\end{Shaded}

\subsection{Tableau stratifié}\label{tableau-stratifiuxe9}

\begin{Shaded}
\begin{Highlighting}[]
\NormalTok{tab\_resume }\OtherTok{\textless{}{-}}\NormalTok{ data\_tab }\SpecialCharTok{\%\textgreater{}\%}
  \FunctionTok{mutate}\NormalTok{(}
    \AttributeTok{sector =} \FunctionTok{as\_factor}\NormalTok{(sector),}
    \AttributeTok{s1q2 =} \FunctionTok{as\_factor}\NormalTok{(s1q2)}
\NormalTok{  ) }\SpecialCharTok{\%\textgreater{}\%}
  \FunctionTok{select}\NormalTok{(sector, s1q4, s1q2, taille\_men) }\SpecialCharTok{\%\textgreater{}\%}
  \FunctionTok{tbl\_summary}\NormalTok{(}
    \AttributeTok{by =}\NormalTok{ sector,}
    \AttributeTok{label =} \FunctionTok{list}\NormalTok{(}
\NormalTok{      s1q4 }\SpecialCharTok{\textasciitilde{}} \StringTok{"Âge"}\NormalTok{,}
\NormalTok{      s1q2 }\SpecialCharTok{\textasciitilde{}} \StringTok{"Sexe"}\NormalTok{,}
\NormalTok{      taille\_men }\SpecialCharTok{\textasciitilde{}} \StringTok{"Taille du ménage"}
\NormalTok{    ),}
    \AttributeTok{type =} \FunctionTok{list}\NormalTok{(}
\NormalTok{      s1q4 }\SpecialCharTok{\textasciitilde{}} \StringTok{"continuous2"}\NormalTok{,}
\NormalTok{      taille\_men }\SpecialCharTok{\textasciitilde{}} \StringTok{"continuous2"}
\NormalTok{    ),}
    \AttributeTok{statistic =} \FunctionTok{list}\NormalTok{(}
      \FunctionTok{all\_continuous}\NormalTok{() }\SpecialCharTok{\textasciitilde{}} \FunctionTok{c}\NormalTok{(}\StringTok{"\{mean\} (\{sd\})"}\NormalTok{, }\StringTok{"\{median\} [\{p25\} ; \{p75\}]"}\NormalTok{),}
      \FunctionTok{all\_categorical}\NormalTok{() }\SpecialCharTok{\textasciitilde{}} \StringTok{"\{n\} (\{p\}\%)"}
\NormalTok{    ),}
    \AttributeTok{missing =} \StringTok{"no"}
\NormalTok{  ) }\SpecialCharTok{\%\textgreater{}\%}
  \FunctionTok{add\_p}\NormalTok{(}
    \AttributeTok{test =} \FunctionTok{list}\NormalTok{(}
\NormalTok{      s1q4 }\SpecialCharTok{\textasciitilde{}} \StringTok{"wilcox.test"}\NormalTok{,}
\NormalTok{      s1q2 }\SpecialCharTok{\textasciitilde{}} \StringTok{"chisq.test"}\NormalTok{,}
\NormalTok{      taille\_men }\SpecialCharTok{\textasciitilde{}} \StringTok{"wilcox.test"}
\NormalTok{    )}
\NormalTok{  ) }\SpecialCharTok{\%\textgreater{}\%}
  \FunctionTok{bold\_labels}\NormalTok{()}

\NormalTok{tab\_resume}
\end{Highlighting}
\end{Shaded}

\begin{table}[t]
\fontsize{12.0pt}{14.0pt}\selectfont
\begin{tabular*}{\linewidth}{@{\extracolsep{\fill}}lccc}
\toprule
\textbf{Characteristic} & \textbf{1. Urban}  N = 8,719\textsuperscript{\textit{1}} & \textbf{2. Rural}  N = 21,618\textsuperscript{\textit{1}} & \textbf{p-value}\textsuperscript{\textit{2}} \\ 
\midrule\addlinespace[2.5pt]
{\bfseries Âge} &  &  & <0.001 \\ 
    Mean (SD) & 25 (19) & 24 (20) &  \\ 
    Median [Q1 ; Q3] & 19 [9 ; 38] & 17 [8 ; 35] &  \\ 
{\bfseries Sexe} &  &  & 0.5 \\ 
    1. MALE & 3,730 (49\%) & 9,459 (50\%) &  \\ 
    2. FEMALE & 3,837 (51\%) & 9,531 (50\%) &  \\ 
{\bfseries Taille du ménage} &  &  & <0.001 \\ 
    Mean (SD) & 7 (4) & 9 (4) &  \\ 
    Median [Q1 ; Q3] & 7 [5 ; 9] & 8 [6 ; 11] &  \\ 
\bottomrule
\end{tabular*}
\begin{minipage}{\linewidth}
\vspace{.05em}
\textsuperscript{\textit{1}}n (\%)\\
\textsuperscript{\textit{2}}Wilcoxon rank sum test; Pearson's Chi-squared test\\
\end{minipage}
\end{table}

Le tableau ci-dessus nous donne la répartition selon le milieu de
résidence de l'âge des individus ainsi que de la taille du ménage; pour
chaque croisement, un test est effectué pour voir la significativité de
la différence de moyenne entre ces deux groupes. On constate que dans le
milieu urbain, l'âge moyen des individus est de 25 ans tandis que dans
le milieu rural, il est de 24 ans. De plus, le test affiche une p-value
très inférieure à 5\%, raison pour laquelle, on peut conclure que la
différence de moyenne entre ces deux groupes est significatif. Pour la
taille des ménages, on constate que la moyenne en milieu rural est de 7
alors qu'en milieu urbain elle est de 9. Le test affiche aussi une
p-value très inférieure à 5\%, ce qui nous permet de conclure que la
différence de moyenne entre ces deux groupes est significatif. Le
tableau présente aussi la répartition selon le sexe et le milieu de
résidence des individus. On constate que 49\% des individus sont de sexe
masculin en milieu urbain et qu'en milieu rural 50\% sont de sexe
masculin. La répartition selon le sexe est donc relativement équilibrée,
ce qui est prouvé par un test de proportion affichant une p-value de
0,5.

\end{document}
